\documentclass[11pt]{article}
\input{../../shared/project_style.tex}
\rhead{Basics 31 Textbook Draft}

\begin{document}
\DocTitle{Basics 31: Introduction to Waves in Plasmas}{V - Plasma waves and Alfvén physics}{Phase, dispersion, polarization, eigenmodes, damping, instability, and plasma wave families}

\begin{overviewbox}
\textbf{Textbook draft, version 1.0.} Plasma waves couple particle inertia, pressure, and electromagnetic restoring forces. This chapter develops a common language for waves before specializing to Alfvén physics: phase and group velocity, polarization, dispersion matrices, continua, damping, and instability. This chapter is designed for a reader with no prior plasma-physics training. It distinguishes exact identities from physical approximations and connects the central equations to later simulation modules.
\end{overviewbox}

\ProjectTOC

\section{Learning objectives}
After completing this note, the reader should be able to:
\begin{itemize}[leftmargin=1.6em]
\item Identify the ingredients required for wave motion.
\item Interpret plane-wave phase and wave vector.
\item Derive a dispersion relation from a linear system.
\item Distinguish phase velocity, group velocity, and polarization.
\item Interpret real and imaginary frequency components.
\item Classify major electrostatic and electromagnetic plasma-wave families.
\end{itemize}

\section{Prerequisites and reading strategy}
\begin{itemize}[leftmargin=1.6em]
\item Basics 5-8, 12, 18, 19, and 29.
\end{itemize}
On a first reading, emphasize physical meaning and dimensional checks. On a second reading, reproduce the derivations. Attempt the computational laboratory only after the limiting cases are understood.

\section{What makes a wave?}

A wave requires a perturbable state, inertia or delayed response, and a restoring mechanism. A schematic field equation,
\begin{equation}
\partial_t^2\psi+\mathcal L\psi=0,
\end{equation}
oscillates when the spatial operator $\mathcal L$ has positive eigenvalues. In plasmas, restoring forces may arise from charge separation, pressure, magnetic tension, or gyromotion, while several species provide distinct inertias.


\section{Plane-wave phase}

A local plane wave is written
\begin{equation}
\psi(\mathbf x,t)=\Re\left\{\tilde\psi
e^{i(\mathbf k\cdot\mathbf x-\omega t)}\right\}.
\end{equation}
Surfaces of constant phase satisfy
\begin{equation}
\mathbf k\cdot\mathbf x-\omega t=\text{constant}.
\end{equation}
$\mathbf k$ points normal to phase fronts and has units of inverse length. The angular frequency $\omega$ has units of inverse time. Sign conventions must be recorded because they determine how growth rates and propagation directions are interpreted.


\section{Dispersion relations from linear equations}

After linearization and Fourier substitution,
\begin{equation}
\partial_t\rightarrow-i\omega,\qquad
\nabla\rightarrow i\mathbf k,
\end{equation}
a coupled PDE system becomes
\begin{equation}
\mathsf D(\omega,\mathbf k)\tilde{\mathbf q}=0.
\end{equation}
A nontrivial amplitude requires
\begin{equation}
\boxed{\det\mathsf D(\omega,\mathbf k)=0}.
\end{equation}
This is the dispersion relation. Multiple roots describe distinct wave branches; the associated null vectors determine polarization and relative phase among field components.


\section{Phase and group velocity}

For a one-dimensional branch $\omega(k)$,
\begin{equation}
v_{\rm ph}=\frac{\omega}{k},\qquad
v_g=\frac{d\omega}{dk}.
\end{equation}
Phase velocity tracks constant phase. Group velocity gives the leading propagation speed of a narrow wave packet when the medium varies slowly and dissipation is weak. In anisotropic plasma, $\mathbf v_g=\nabla_{\mathbf k}\omega$ need not be parallel to $\mathbf k$.


\section{Polarization}

Polarization describes the direction and phase relation of vector perturbations. Examples include:
\begin{itemize}[leftmargin=1.6em]
\item longitudinal electrostatic oscillations with $\tilde{\mathbf E}\parallel\mathbf k$;
\item transverse electromagnetic waves with $\tilde{\mathbf E}\perp\mathbf k$;
\item circularly polarized waves in magnetized plasma;
\item shear-Alfvén waves with transverse velocity and magnetic perturbations.
\end{itemize}
A scalar dispersion curve alone does not specify the physical mode.


\section{Standing waves, eigenmodes, and continua}

Boundaries or spatially varying coefficients replace a freely chosen $k$ by a global eigenproblem:
\begin{equation}
\mathcal A\boldsymbol\xi=\omega^2\mathcal M\boldsymbol\xi.
\end{equation}
A finite domain often produces discrete standing modes. In nonuniform ideal MHD, each magnetic surface can carry its own local Alfvén frequency, generating a continuum. Geometry-induced coupling may open gaps in that continuum where discrete global eigenmodes can exist.


\section{Complex frequency, damping, and growth}

Using $e^{-i\omega t}$, let
\begin{equation}
\omega=\omega_r+i\gamma.
\end{equation}
Then
\begin{equation}
e^{-i\omega t}=e^{-i\omega_rt}e^{\gamma t}.
\end{equation}
Thus $\gamma>0$ denotes growth and $\gamma<0$ damping under this convention. Mechanisms include collisions, phase mixing, continuum damping, Landau damping, and energetic-particle drive. A stable oscillation can still decay because of nonideal or kinetic processes.


\section{Representative plasma-wave families}

At elementary level:
\begin{description}[leftmargin=2.5em]
\item[Langmuir waves:] electron charge-separation oscillations near $\omega_{pe}$.
\item[Ion-acoustic waves:] compressive electrostatic waves involving electron pressure and ion inertia.
\item[Electromagnetic waves:] light-like modes modified by plasma response.
\item[Alfvén waves:] magnetic-tension waves propagating primarily along the field.
\item[Magnetosonic waves:] compressive coupled pressure-magnetic waves.
\end{description}
The appropriate model depends on frequency ordering, propagation angle, collisionality, and scale.


\section{Computational laboratory}
Solve several analytic dispersion relations and plot $\omega(k)$, phase velocity, and group velocity. Construct a two-component dispersion matrix and recover both eigenvalues and polarization vectors. Form a wave packet numerically and compare measured envelope speed with $d\omega/dk$.

\section{Summary}
\begin{itemize}[leftmargin=1.6em]
\item Plasma waves combine inertia with electrostatic, pressure, magnetic, or gyroscopic restoring forces.
\item Dispersion relations determine allowed frequencies; eigenvectors determine polarization.
\item Phase and group velocities describe different aspects of propagation.
\item Nonuniform plasmas can support continua and discrete global eigenmodes.
\item Complex frequency unifies oscillation, damping, and instability under a stated sign convention.
\end{itemize}

\SymbolsAcronymsSection
\begin{symtable}
$\mathbf k$ & wave vector & Spatial phase gradient. \\
$\omega$ & complex angular frequency & Oscillation and growth/damping rate. \\
$v_{\rm ph}$ & phase velocity & $\omega/k$ in one dimension. \\
$v_g$ & group velocity & $d\omega/dk$. \\
$\mathsf D$ & dispersion matrix & Fourier-transformed linear operator. \\
$\gamma$ & growth rate & $\operatorname{Im}\omega$ for $e^{-i\omega t}$. \\
\end{symtable}

\section{Exercises}
\begin{enumerate}[leftmargin=1.8em]
\item Derive the constant-phase propagation velocity of a plane wave.
\item Obtain the wave-equation dispersion relation by Fourier substitution.
\item Explain why $\det\mathsf D=0$ is required for a nonzero amplitude.
\item For a given dispersive branch, calculate phase and group velocities.
\item Translate the growth-rate sign if the time convention is changed to $e^{+i\omega t}$.
\end{enumerate}

\section{References}
\begin{thebibliography}{99}
\bibitem{ref1} G. B. Arfken, H. J. Weber, and F. E. Harris, \emph{Mathematical Methods for Physicists}, 7th ed., Academic Press (2013).
\bibitem{ref2} G. Strang, \emph{Linear Algebra and Its Applications}, 4th ed., Brooks/Cole (2006).
\bibitem{ref3} W. A. Strauss, \emph{Partial Differential Equations: An Introduction}, 2nd ed., Wiley (2007).
\bibitem{ref4} F. F. Chen, \emph{Introduction to Plasma Physics and Controlled Fusion}, 3rd ed., Springer (2016).
\bibitem{ref5} R. J. Goldston and P. H. Rutherford, \emph{Introduction to Plasma Physics}, CRC Press (1995).
\bibitem{ref6} P. M. Bellan, \emph{Fundamentals of Plasma Physics}, Cambridge University Press (2006).
\bibitem{ref7} T. H. Stix, \emph{Waves in Plasmas}, American Institute of Physics (1992).
\bibitem{ref8} D. A. Gurnett and A. Bhattacharjee, \emph{Introduction to Plasma Physics}, 2nd ed., Cambridge University Press (2017).
\bibitem{ref9} J. P. Freidberg, \emph{Ideal MHD}, Cambridge University Press (2014).
\bibitem{ref10} J. P. Goedbloed, R. Keppens, and S. Poedts, \emph{Advanced Magnetohydrodynamics}, Cambridge University Press (2010).
\bibitem{ref11} H. Goedbloed and S. Poedts, \emph{Principles of Magnetohydrodynamics}, Cambridge University Press (2004).
\end{thebibliography}

\end{document}
